\chapter{Lecture 11: Copy-and-Swap, Rvalues and Move Semantics}
\begin{center}
\textit{February 8, 2018}
\end{center}

\section{Recap and the Failing-\code{new} Problem}
\begin{lstlisting}[style=cpp]
Node &Node::operator=(const Node &other) {
    if (this == &other) return *this;
    data = other.data;
    delete next;
    next = other.next ? new Node{*other.next} : nullptr;
    return *this;
}
\end{lstlisting}

\paragraph{Q.} What if \code{new} fails?

\code{new} throws an exception and the function is exited before \code{next} is
assigned. If the caller catches the exception and tries to recover, the object's
\code{next} is now a dangling pointer (and our dtor will delete it again later).

\paragraph{Solution.} Don't delete the old data until the new data is safely in
hand:
\begin{lstlisting}[style=cpp]
Node &Node::operator=(const Node &other) {
    if (this == &other) return *this;
    Node *tmp = next;
    next = other.next ? new Node{*other.next} : nullptr;  // if this throws,
    data = other.data;                                    // nothing has changed
    delete tmp;
    return *this;
}
\end{lstlisting}
(In fact this version is correct even without the self-assignment check, but the
check saves work.)

\section{The Copy-and-Swap Idiom}
\begin{lstlisting}[style=cpp]
// node.h
struct Node {
    ...
    void swap(Node &other);
    Node &operator=(const Node &other);
};

// node.cc
#include <utility>
void Node::swap(Node &other) {
    using std::swap;
    swap(data, other.data);
    swap(next, other.next);
}

Node &Node::operator=(const Node &other) {
    Node tmp = other;   // deep copy via the copy ctor (may throw: fine)
    swap(tmp);          // I take tmp's fresh copy; tmp takes my old data
    return *this;       // tmp's dtor frees my old data
}
\end{lstlisting}
This relies on a correctly implemented copy ctor and dtor, and is safe for
self-assignment and for a failing \code{new} automatically.

\section{Member Functions vs.\ Standalone Functions}
\code{operator=} must be implemented as a method; the left-hand side is
represented by \code{*this}. Other operators can be either.

\begin{example}
Implement the Vec operators as methods.
\begin{lstlisting}[style=cpp]
// vec.h
struct Vec {
    int x, y;
    Vec operator+(const Vec &other) const;
    Vec operator*(int k) const;            // only v * k!
};

// vec.cc
#include "vec.h"
Vec Vec::operator+(const Vec &other) const {
    return {x + other.x, y + other.y};
}
Vec Vec::operator*(int k) const {
    return {x * k, y * k};
}
\end{lstlisting}
\code{k * v} cannot be a method of \code{Vec}, since the left operand is an
\code{int}. It must remain a standalone function:
\code{Vec operator*(int k, const Vec \&v) \{ return v * k; \}}.
\end{example}

What about I/O operators as methods?
\begin{lstlisting}[style=cpp]
struct Vec {
    ...
    std::ostream &operator<<(std::ostream &out) const {
        out << x << " " << y;
        return out;
    }
};

Vec v{1, 2};
v << cout;       // this compiles... and looks completely backwards
\end{lstlisting}
\begin{important}[Rule]
I/O operators should always be implemented as standalone functions, because the
stream must be the left operand.
\end{important}

\section{Tracing Copies}
\begin{lstlisting}[style=cpp]
Node plusOne(Node n) {
    for (Node *p = &n; p; p = p->next) ++p->data;
    return n;
}

Node n{1, new Node{2, nullptr}};
Node n2{plusOne(n)};
\end{lstlisting}
\begin{itemize}
  \item \code{n} is built with the basic ctor.
  \item Passing \code{n} by value to \code{plusOne}: copy ctor (deep copy of both
        nodes).
  \item Returning \code{n} by value from \code{plusOne}: copy into a temporary
        in \code{main}'s frame.
  \item Initializing \code{n2} from that temporary: copy ctor again.
\end{itemize}
Compiling with \code{g++ -std=c++14 -fno-elide-constructors node.cc} shows all
the copy-ctor calls; without the flag the compiler removes some of them (copy
elision, next lecture).

Look at the last step: \code{n2} is copy-constructed from a \emph{temporary},
which is destroyed as soon as \code{n2} has been built. Deep-copying the whole
list, only to throw the original away a moment later, is wasteful. When we are
about to copy from a temporary, it makes sense to not copy but to take over
(steal) its data.

\section{Rvalues and Rvalue References}
\paragraph{Q.} How do we know a value is a temporary?

\begin{definition}
An \textbf{lvalue} is anything with an address (a named storage location). An
\textbf{rvalue} is a value that is not an lvalue: temporaries such as
\code{plusOne(n)}, \code{x + y}, or literals.
\end{definition}

C++11 introduced \textbf{rvalue references}, written \code{T\&\&}, which bind
to rvalues: references to temporary objects that are about to be destroyed.
Overloading on \code{const T\&} vs.\ \code{T\&\&} lets the compiler pick
``copy'' for lvalues and ``steal'' for rvalues.

\section{The Move Constructor}
The move ctor takes a single rvalue reference as its parameter:
\begin{lstlisting}[style=cpp]
Node::Node(Node &&other) : data{other.data}, next{other.next} {
    other.next = nullptr;   // other is about to be destroyed and its dtor
                            // will run: make sure we steal, not share
}
\end{lstlisting}

\begin{center}
\begin{tikzpicture}[box/.style={draw, minimum width=13mm, minimum height=7mm},
  >={Stealth}]
\node[box] (o) at (0,0) {1 \,/};
\node[left=2mm of o] {\code{other}};
\node[box] (t) at (0,-1.5) {1 \,$\bullet$};
\node[left=2mm of t] {\code{this}};
\node[box] (b) at (3,-0.75) {2 \,$\bullet$};
\node[box] (c) at (5.5,-0.75) {3 \,/};
\draw[->] (t) -- (b);
\draw[->] (b) -- (c);
\draw[->, dashed, gray] (o) -- node[above, sloped, font=\scriptsize] {was} (b);
\end{tikzpicture}
\end{center}

\section{The Move Assignment Operator}
\begin{lstlisting}[style=cpp]
Node m{10, new Node{20, nullptr}};
m = plusOne(n);           // the right side is an rvalue: move assignment

Node &Node::operator=(Node &&other) {
    using std::swap;
    swap(data, other.data);
    swap(next, other.next);   // or simply: swap(other);
    return *this;
}
\end{lstlisting}
Swapping gives \code{other} our old list; \code{other} is a temporary, so its
dtor frees our old list for us. No allocation, so it cannot fail.

If you don't write the move operations, the copy versions are used for rvalues
too (it is still correct, just slower).

\section{Copy/Move Elision (Preview)}
\begin{lstlisting}[style=cpp]
// c++/elision/vec.cc
Vec makeAVec() { return {0, 0}; }
Vec v = makeAVec();
\end{lstlisting}
Theoretically: one basic ctor (inside \code{makeAVec}) plus one copy/move ctor
(to get the value into \code{v}). In practice only the basic ctor runs: the
\code{Vec\{0,0\}} is created directly in the space for \code{v}. More next
lecture.

\section*{Summary}
\begin{itemize}
  \item In \code{operator=}, allocate the new data before deleting the old, or
        use copy-and-swap.
  \item Operators whose left operand isn't your class (\code{k * v},
        \code{<<}, \code{>>}) must be standalone functions.
  \item Rvalues are temporaries; \code{T\&\&} binds to them.
  \item Move ctor/assignment steal the resources of a dying object and leave it
        in a destructible state (e.g.\ \code{next = nullptr}).
\end{itemize}
